\section{Cone pictures, asphericity, and protected roots}\label{sec:topology}

The probabilistic argument has produced finite immersed graphs for which
no reduced spherical arrangement exists.  We now turn that exclusion
into the two geometric properties needed in Section~\ref{sec:algebra}:
asphericity of the cone complex and protection of the roots.  We give
the cone-picture argument of
\cite[Proposition~5.1 and Corollary~5.3]{OpenAI2026ZD} in full,
including the surgeries that
preserve a rooted boundary.

\begin{proposition}[Cone criterion]\label{prop:cone-criterion}
Let $\Gamma=\Gamma_A\sqcup\Gamma_B$ be a finite graph immersed in the
rose $F$, with chosen roots $x_A,x_B$, and let $X$ and $G$ be defined
by \eqref{eq:cone-complex}.  If $\Gamma$ admits no reduced spherical
arrangement as in Definition~\ref{def:reduced-arrangement}, then
\[
 \pi_2(X)=0
 \quad\text{and}\quad
 \text{the root-protection condition \eqref{eq:root-protection} holds}.
\]
Moreover, $X$ is a finite aspherical complex of dimension at most two,
and $G$ is finitely presented and torsion-free.
\end{proposition}

The main issue is to pass from a sphere or a null-homotopy to a reduced
arrangement.  Its disks will record fillings through cones, and its arcs
will record inverse letters in the intervening map to the rose.  When a
pair of inverse letters traverses the same edge of $\Gamma$, a band
surgery removes that pair.  Minimizing the total boundary length makes
all such pairs impossible.

\subsection{Maps represented by cone pictures}

For $\Sigma=S^2$ or $D^2$, a \emph{cone picture} is a map
$f\colon\Sigma\to X$ together with finitely many pairwise disjoint
piecewise linear closed disks $D_i$ in the interior of $\Sigma$ with
the following properties.  On the complement of their interiors, $f$
takes values in $F$.  For each $i$ there is a component $\Lambda_i$ of
$\Gamma$ and a factorization
\[
 f|_{D_i}\colon D_i\longrightarrow K_{\Lambda_i}\longrightarrow X
\]
whose restriction to $\partial D_i$ is a specified closed edge path
in $\Lambda_i$.  In the disk case we also specify an edge path in
$\Gamma$ whose label is read on $\partial\Sigma$.  That path may have
distinct endpoints, since all graph vertices map to $*$.

\begin{lemma}[Picture representation]\label{lem:cone-pictures}
Every map $S^2\to X$ is homotopic to a cone picture.  Every disk map
whose boundary is the label of a specified edge path in $\Gamma$ is
homotopic relative to that boundary to a cone picture.
\end{lemma}

\begin{proof}
We first give a relative replacement of each cone by ordinary relator
disks.  Choose a maximal tree in a component $\Lambda$ and a vertex of
that tree.  The edges outside the tree give a free basis of
$\pi_1(\Lambda)$, represented by based edge loops.  Let $Y_\Lambda$
be $\Lambda$ with one disk attached along each of these basis loops.
Van Kampen's theorem gives $\pi_1(Y_\Lambda)=0$.  The cellular boundary
map from its two-cells has image the integral cycle group of $\Lambda$
and is injective: the basis loops give a basis of
$H_1(\Lambda;\mathbb Z)$.  Therefore $Y_\Lambda$ has zero reduced
integral homology.  A simply connected acyclic CW complex is
contractible, by the Hurewicz and Whitehead theorems
\cite[Sections 4.1 and 4.2]{Hatcher2002}.

Both $Y_\Lambda$ and $K_\Lambda$ contain $\Lambda$ as a CW subcomplex
and are contractible.  The identity on $\Lambda$ extends to maps in
both directions: extend successively over the relative cells, using
the vanishing homotopy groups of the target.  The two composites are
homotopic to the corresponding identities relative to $\Lambda$.
For this last assertion use the same cell-extension argument on the
products with $[0,1]$, prescribing the homotopy on the two ends and on
$\Lambda\times[0,1]$.  Thus the replacements are homotopy equivalences
relative to $\Lambda$, not merely equivalences of their absolute
homotopy types.  After attachment along $\Lambda\to F$, they give a
homotopy equivalence relative to $F$ between $X$ and the finite
presentation complex obtained from $F$ by attaching these relator
disks.

Here is the picture construction for this presentation complex.
Replace the attaching maps by mapping-cylinder collars and subdivide,
so that simplicial approximation applies to the surface map.  This
changes neither the homotopy type relative to $F$ nor a prescribed
boundary word.  In the interior of each relator disk choose a point
outside the images of the one-dimensional simplices and of the
degenerate two-dimensional simplices of the surface.  Its preimage is
a finite set of points, each with a neighborhood mapped
homeomorphically to a neighborhood of the chosen point.  By compactness
and finiteness of the triangulations, a sufficiently small closed disk
about the chosen point has as its entire preimage a disjoint union of
such disk neighborhoods.

In each target relator disk, expand that small disk to the full disk
and collapse the surrounding annulus radially to its boundary.  This
map is homotopic to the identity relative to the boundary.  Collapse
the mapping-cylinder collars to their attaching paths in $F$ as well.
After postcomposition, the surface outside the selected source disks
maps into $F$, while each source disk maps into one relator disk.
Returning through the relative equivalence to $X$ makes each of these
disk maps factor through the corresponding abstract cone.  All
homotopies fix a prescribed outer boundary.  This gives the asserted
cone pictures.
\end{proof}

We shall also use two elementary changes to a picture.  A homotopy of
one inner boundary path in its graph can be made on a collar of that
boundary.  Fill the resulting boundary in the same cone.  The new and
old fillings, with the collar included, are homotopic relative to the
original boundary because the abstract cone is contractible.  Thus
inner paths can be tightened cyclically.  If the result is constant,
the inner disk can be removed, filling it by that constant map into
$F$.  An outer path can likewise be tightened with its graph endpoints
fixed; the corresponding homotopy of its label is accommodated by a
boundary collar.

\subsection{Length reduction by band surgery}

Suppose first that $\pi_2(X)\ne0$.  Among cone pictures representing
non-null-homotopic spheres, choose one minimizing the sum of the lengths
of its inner boundary paths.  Alternatively, suppose root protection
fails.  Fix a root $r\in\{x_A,x_B\}$ and a distinct vertex $z$ joined
to it by a path with trivial label in $G$.  Among all cone pictures of
null-homotopies of labels of paths from $r$ to this fixed $z$, choose
one minimizing the sum of all inner boundary lengths and the outer
path length.  These minima exist because lengths are nonnegative
integers and Lemma~\ref{lem:cone-pictures} provides the required
pictures.  Among pictures attaining the minimum, remove all constant
inner boundaries.

The preceding collar changes show that every remaining inner boundary
is nonempty and cyclically immersed.  In the disk case the outer path
is immersed and nonempty, since its endpoints are distinct.  There is
at least one inner disk.  For the sphere case, a map with no such disk
would factor through $F$, whose universal cover is a tree and whose
second homotopy group is zero.  For the disk case, a nonempty immersed
path in $\Gamma$ has a nonempty freely reduced label word, which
cannot bound a disk in $F$.

Complete a disk picture to a sphere by adjoining a formal exterior
disk $D_\infty$.  Its interior carries no map into $X$.  Mark its
boundary at the initial point of the specified outer path; this is
the exceptional break, representing the two endpoints of that path.
Choose the orientation of the sphere so that the boundary of
$D_\infty$, starting at the break, reads the path from $r$ to $z$.
This choice may reverse all ordinary boundary words, which preserves
cyclic immersion.  In the sphere case simply orient the sphere and
give all disks their induced orientations.

On the complement of these disks, the map takes values in the rose.
Put it in piecewise linear general position with respect to one point
in the interior of each rose edge, keeping the boundary words fixed.
The preimages of these points are disjoint embedded arcs and circles.
Every letter occurrence has exactly one arc endpoint.  Discard the
circles.  The arcs pair all occurrences, and paired letters are inverse
when read in the oriented disk boundaries: the two ends of a
transverse arc have opposite crossing signs, and disk-boundary
orientations are opposite to those of the complementary surface.

It remains to show that paired occurrences use distinct underlying
edges of $\Gamma$.  Suppose an arc pairs two traversals of the same
edge $e$.  Choose a thin rectangular band around the arc.  It maps
into a small open interval in the corresponding rose edge and lifts
to the interior of the abstract graph edge $e$.  The lift agrees with
both specified boundary lifts at the band ends.  Whenever an inner
disk meets the band, its cone must therefore be the cone on the
component containing $e$.

The local shortening is particularly simple in graph paths.  Splice
the disk boundaries along the two sides of the band and then contract
the residual portions of the two paired traversals within $e$.
This removes two full edge occurrences in total. For ordinary
boundaries the changes, with suitable cyclic starting points, are
\begin{equation}\label{eq:inner-surgery-words}
 (eP,\ e^{-1}Q)\longmapsto PQ,
 \qquad
 eP e^{-1}Q\longmapsto(P,Q).
\end{equation}
Here $e^{-1}$ denotes the reverse oriented graph edge, rather than
merely a letter with the same label.  Consequently all concatenations
in \eqref{eq:inner-surgery-words} are paths in the same graph, and the
new boundaries are closed.  For a rooted exterior boundary the
corresponding changes are
\begin{equation}\label{eq:exterior-surgery-words}
 (P e R,\ e^{-1}Q)\longmapsto PQR,
 \qquad
 P e Q e^{-1}R\longmapsto PR.
\end{equation}
In the second formula $Q$ is closed at the terminal vertex of $e$;
in both formulas the resulting exterior path has exactly the original
graph endpoints $r,z$.  The contractions along $e$ are realized by
boundary-collar homotopies.

We now specify the resulting surface maps in all four cases.  This
also explains why the operation retains the failure that was
minimized.

If the band joins two distinct inner disks, their union with the band
is a disk.  The maps on all three pieces factor together through the
same abstract cone, since their lifts agree on the attaching intervals.
Replace them by this single cone disk and shorten its boundary as in
the first formula of \eqref{eq:inner-surgery-words}.  The replacement
is homotopic to the old map relative to the rest of the surface, and
the total boundary length decreases by two.  Further tightening can
only decrease it.  An essential sphere remains essential, and a disk
still fills the same exterior word.

If the band joins an inner disk to $D_\infty$, their union is again a
disk, now an enlarged exterior disk.  Retain its complementary disk
and restrict the old map to that complement.  Its new boundary is
$PQR$ as in the first formula of \eqref{eq:exterior-surgery-words}.
The resulting picture is a null-homotopy of a path from $r$ to $z$;
the old inner boundary and exterior boundary have been replaced by a
boundary shorter than their combined length by two.

If both ends of the band lie on $D_\infty$, their union is an annulus
in the completed sphere.  Its complement consists of two disks.
Retain the one whose boundary contains the marked break, as shown
in Figure~\ref{fig:exterior-surgery}.  The restricted map on that disk,
after the edge-collar contraction, fills the word $PR$ in the second
formula of \eqref{eq:exterior-surgery-words}.  Thus this operation uses
only a restriction of the existing map; it does not require a
separate filling of $Q$.  It preserves the fixed endpoints and
decreases the total length by at least two.  The path $PR$ cannot
tighten to the empty path, because $r\ne z$.

\begin{figure}[H]
\centering
\begin{tikzpicture}[x=1cm,y=1cm,font=\small,
  line cap=round,line join=round,>=Latex]
  \node[anchor=west,font=\small\bfseries] at (0,3.00)
    {Band returning to the exterior disk};
  \fill[black!10] (0.15,-1.70) rectangle (5.00,2.45);
  \fill[white] (2.575,0.35) ellipse (2.10cm and 1.55cm);
  \fill[black!20] (0.475,0.19) rectangle (4.675,0.51);
  \draw (2.575,0.35) ellipse (2.10cm and 1.55cm);
  \draw (0.49,0.19) -- (4.66,0.19);
  \draw (0.49,0.51) -- (4.66,0.51);
  \node[font=\footnotesize] at (2.575,2.17)
    {formal exterior disk $D_\infty$};
  \node[font=\footnotesize,fill=black!20,inner sep=1pt]
    at (2.575,0.35) {band};
  \node at (2.575,1.40) {$Q$};
  \node[font=\footnotesize] at (2.575,0.94) {$Q$ side};
  \node at (1.24,-0.75) {$P$};
  \node at (3.91,-0.75) {$R$};
  \node[font=\footnotesize] at (2.575,-0.45) {retain this side};
  \draw[->,line width=0.9pt] (0.52,-0.01) -- (0.52,0.73);
  \draw[->,line width=0.9pt] (4.63,0.73) -- (4.63,-0.01);
  \node[anchor=east] at (0.37,0.35) {$e$};
  \node[anchor=west] at (4.80,0.35) {$e^{-1}$};
  \fill (2.575,-1.20) circle (1.7pt);
  \node[font=\footnotesize,align=center] at (2.575,-1.51)
    {break: fixed graph endpoints};

  \draw[->,line width=0.8pt] (5.55,0.35) -- (6.72,0.35);
  \node[font=\footnotesize,align=center] at (6.14,1.03)
    {retain the\\break side};

  \draw (9.15,0.35) ellipse (1.72cm and 1.23cm);
  \node[font=\footnotesize,align=center] at (9.15,0.54)
    {retained disk\\outer path $PR$};
  \node at (8.10,-0.45) {$P$};
  \node at (10.20,-0.45) {$R$};
  \fill (9.15,-0.88) circle (1.7pt);
  \node[font=\footnotesize] at (9.15,-1.20) {same endpoints};
  \node[font=\footnotesize,align=center] at (11.95,1.61)
    {$Q$ boundary\\discarded};
  \draw[densely dashed,black!65] (11.95,0.47) circle (0.61cm);
  \node at (11.95,0.47) {$Q$};
  \node[font=\footnotesize] at (6.55,-2.04)
    {$P\,e\,Q\,e^{-1}\,R\quad\longmapsto\quad PR$};

\end{tikzpicture}
\caption{A band returning to the formal exterior disk, shown
schematically in the domain sphere.  The two paired occurrences
traverse the same graph edge in opposite directions.  The exterior
disk and the band form an annulus.  Retaining its complementary disk
on the break side replaces $PeQe^{-1}R$ by $PR$ with the same graph
endpoints: $Q$ is closed at the terminal vertex of $e$, and the
restricted picture fills $PR$.  Shading denotes the formal exterior
disk and band; the exterior disk itself carries no map to $X$.}
\label{fig:exterior-surgery}
\end{figure}


Finally suppose both band ends lie on one inner disk.  The union $N$
of the disk and band is an annulus, and its map factors through the
same abstract cone $K_\Lambda$.  The complement of $N$ in the sphere
consists of two disks $Q_1,Q_2$.  Cap each new boundary by a disk
mapping through $K_\Lambda$, using the boundary lifts inherited from
$N$, and shorten these boundary paths to $P$ and $Q$ in the second
formula of \eqref{eq:inner-surgery-words}.  In the rooted case one of
$Q_1,Q_2$ contains $D_\infty$.  Retain that component with its new
cap and remove the formal exterior disk.  This is another disk
picture with the same exterior path.  Its total inner boundary length
is at most the old total minus two, because the two possible new cap
boundaries together have length two less than the old boundary and
all other retained disks already belonged to the old picture.

In the sphere case, the same operation gives two capped spheres; at
least one must still be essential.  To verify this without an
ambiguity from choices of caps, lift the original sphere map to the
universal cover $\widetilde X$.  The factorization $N\to K_\Lambda\to
X$ lifts through a single lift $K_\Lambda\to\widetilde X$, chosen to
agree with the lifted sphere on $N$.  Use that very same lift for
both new caps.  Write $\widetilde f_1,\widetilde f_2$ for the resulting
lifted capped spheres, with the orientations induced from $Q_1,Q_2$.
Their fundamental homology classes satisfy
\begin{equation}\label{eq:sphere-surgery-homology}
 \widetilde f_*[S^2]
   =\widetilde f_{1*}[S^2]+\widetilde f_{2*}[S^2]
       \quad\text{in }H_2(\widetilde X;\mathbb Z).
\end{equation}
Compute this identity before the boundary-collar shortenings, which
are homotopies and do not change these homology classes.  The
difference of the cycles on the two sides consists of the
old annulus and the two caps, with their boundary orientations
canceling.  This cycle factors through the contractible abstract cone
$K_\Lambda$, so its homology class is zero.  No assertion that the
image of this cone is an embedded contractible subspace is needed.
Since $\widetilde X$ is simply connected, the Hurewicz map
$\pi_2(\widetilde X)\to H_2(\widetilde X;\mathbb Z)$ is an
isomorphism.  The left side of \eqref{eq:sphere-surgery-homology} is
nonzero, and hence one capped sphere is essential.  Its total boundary
length is smaller than the original one by at least two.

Figure~\ref{fig:inner-surgery} shows this splitting into two capped
spheres and the homology identity that preserves essentiality.

\begin{figure}[H]
\centering
\begin{tikzpicture}[x=1cm,y=1cm,font=\small,
  line cap=round,line join=round,>=Latex]
  \node[anchor=west,font=\small\bfseries] at (0,-2.82)
    {Band returning to an inner cone disk};
  \fill[black!18,even odd rule]
    (2.20,-4.62) circle (1.19cm)
    (2.20,-4.62) circle (0.68cm);
  \draw (2.20,-4.62) circle (1.19cm);
  \draw (2.20,-4.62) circle (0.68cm);
  \node at (2.20,-4.62) {$D_Q$};
  \node at (2.20,-3.15) {$D_P$};
  \node[font=\footnotesize,fill=white,inner sep=1pt]
    at (2.20,-3.68) {annulus};
  \node[anchor=west] at (3.39,-4.62) {$P$};
  \node[anchor=west] at (2.85,-4.62) {$Q$};
  \node[font=\footnotesize,align=center] at (2.20,-6.18)
    {$ePe^{-1}Q$ splits into $P,Q$};

  \draw[->,line width=0.8pt] (4.18,-4.62) -- (5.29,-4.62);
  \node[font=\footnotesize,align=center] at (4.75,-3.88)
    {cap\\both sides};

  \draw (6.70,-4.62) circle (1.05cm);
  \draw (10.18,-4.62) circle (1.05cm);
  \draw[fill=black!18] (6.70,-5.05) ellipse (0.70cm and 0.40cm);
  \draw[fill=black!18] (10.18,-5.05) ellipse (0.70cm and 0.40cm);
  \node at (6.70,-4.25) {$D_P$};
  \node at (10.18,-4.25) {$D_Q$};
  \node[font=\footnotesize] at (6.70,-5.05) {cap};
  \node[font=\footnotesize] at (10.18,-5.05) {cap};
  \node[font=\footnotesize,fill=white,inner sep=0.5pt]
    at (6.25,-4.73) {$P$};
  \node[font=\footnotesize,fill=white,inner sep=0.5pt]
    at (9.73,-4.73) {$Q$};
  \node at (8.43,-4.62) {$+$};
  \node[font=\footnotesize] at (8.44,-3.23)
    {two capped spheres};
  \node[font=\footnotesize] at (8.44,-6.05)
    {Both caps use the same lift of $K_\Lambda\to X$.};
  \node at (8.44,-6.60)
    {$[\widetilde f]=[\widetilde f_1]+[\widetilde f_2]$};
\end{tikzpicture}
\caption{An inner cone disk and its returning band form an annulus,
shown schematically in the domain sphere.  Here $D_P,D_Q$ are its
complementary disks, with boundary paths $P,Q$.  Capping both through
the same abstract cone produces two spheres.  For the original sphere
map $f$, the symbols $[\widetilde f]$ and $[\widetilde f_i]$ denote the
second homology classes of its lift and of the two capped spheres in
the universal cover.  Using for both caps the same lift of
$K_\Lambda\to X$ as for the original annulus gives the displayed
identity.  If $f$ is essential, at least one
capped sphere remains essential.  Shading denotes the annulus and
caps in the domain, without asserting that their images embed in $X$.}
\label{fig:inner-surgery}
\end{figure}


Each case contradicts minimality.  The chosen pairing therefore uses
distinct underlying graph edges.  All the other requirements of
Definition~\ref{def:reduced-arrangement} have already been verified:
the ordinary boundaries are positive and cyclically immersed, there
is at least one such boundary, and the only possible exceptional
boundary is a positive immersed path starting at a designated root.
Thus either failure would give a reduced spherical arrangement.  Its
assumed absence proves both $\pi_2(X)=0$ and root protection.

\subsection{The finite aspherical complex and torsion}

For completeness, we finish the two standard consequences needed in
Proposition~\ref{prop:cone-criterion}.  The universal cover
$\widetilde X$ is simply connected, and the vanishing of $\pi_2(X)$
gives $H_2(\widetilde X;\mathbb Z)=0$ by Hurewicz.  Its higher cellular
homology groups vanish because its dimension is at most two.  Thus it
is simply connected and acyclic.  If it had a nonzero homotopy group,
the first such group would, by Hurewicz, give a nonzero homology group.
All homotopy groups therefore vanish, and Whitehead's theorem makes
$\widetilde X$ contractible.  Hence $X$ is aspherical.  Its finite
two-dimensional CW structure gives a finite presentation for $G$.

The cellular chains of this contractible universal cover give a free
resolution over $\mathbb ZG$,
\[
 0\longrightarrow C_2(\widetilde X)
 \longrightarrow C_1(\widetilde X)
 \longrightarrow C_0(\widetilde X)
 \longrightarrow\mathbb Z\longrightarrow0.
\]
If $G$ had a nonidentity element of finite order, a power of that
element would generate a cyclic subgroup $C_\ell$ of prime order
$\ell$.  Restricting this resolution to $\mathbb ZC_\ell$ still gives
a free resolution of length at most two: $\mathbb ZG$ is a free left
$\mathbb ZC_\ell$-module, with basis given by representatives for the
cosets $C_\ell g$ in $G$.  It would follow that
$H^k(C_\ell;\mathbb F_\ell)=0$ for every $k>2$.

To see the contradiction directly, write $u$ for a generator of
$C_\ell$ and $N=1+u+\cdots+u^{\ell-1}$.  The periodic free resolution
\[
 \cdots\xrightarrow{\,N\,}\mathbb ZC_\ell
 \xrightarrow{\,u-1\,}\mathbb ZC_\ell
 \xrightarrow{\,N\,}\mathbb ZC_\ell
 \xrightarrow{\,u-1\,}\mathbb ZC_\ell
 \longrightarrow\mathbb Z\longrightarrow0
\]
is exact: the augmentation ideal is the image of multiplication by
$u-1$, the kernel of that multiplication is $\mathbb ZN$, and the
kernel of multiplication by $N$ is the
augmentation ideal.  Applying
$\operatorname{Hom}_{\mathbb ZC_\ell}(-,\mathbb F_\ell)$ with trivial
coefficients makes every differential zero, since $u-1$ acts as zero
and $N$ acts as $\ell=0$.  Therefore
$H^k(C_\ell;\mathbb F_\ell)\cong\mathbb F_\ell$ for all $k\ge0$,
the standard cohomological obstruction to torsion in a group of finite
cohomological dimension \cite[Chapter VIII]{Brown1982}.  This
contradiction proves that $G$ is torsion-free and completes the proof
of Proposition~\ref{prop:cone-criterion}.
